\chapter{Gaussian Processes and Kriging}

\vfill
\localtableofcontents
\vfill
\newpage\label{sec:GP}

% ============================================================
\section{Introduction}
% ============================================================

The analytical and Monte-Carlo methods presented in the previous chapters
assume that the observation function $h_k$ is fully known. In many practical
navigation scenarios, however, $h_k$ depends on a physical field --- such as
a terrain elevation map or a magnetic anomaly map --- that is either unavailable
or known only through scarce and noisy measurements. Gaussian processes (GPs)
provide a principled probabilistic framework for modelling such functions from
data: a GP places a prior distribution directly over functions, encoding both a
prior expectation of the field and its spatial correlation structure through a
mean function and a covariance kernel.

This chapter develops the GP framework and three distinct ways in which it is
exploited in this thesis. First, GP regression (kriging) yields the optimal
posterior estimate of the field at any query point given noisy training data;
the resulting estimator and its uncertainty are derived in closed form, and serve
as the observation model for the particle filters described in later chapters.
Second, the spectral structure of the covariance kernel underlies reduced-rank
approximations, in which the GP is projected onto a finite set of basis
functions; these approximations make the otherwise cubic cost of exact GP
inference tractable in high-dimensional or large-data settings. Third, every GP
covariance kernel canonically induces a reproducing kernel Hilbert space (RKHS),
which provides a deterministic, function-analytic interpretation of GP
regression: the posterior mean is the minimum-norm interpolant in the RKHS, and
the kernel encodes the smoothness class of the unknown field. This RKHS
perspective is the foundation of kernel Kalman filtering methods, introduced in
Section~\ref{sec:kernel_kalman}.

The structure of the chapter follows these three threads. Section~\ref{sec:GP_definition}
introduces GPs and their covariance kernels, including the squared exponential
and Mat\'{e}rn families. Section~\ref{sec:hyperparameters} addresses hyperparameter
estimation by marginal likelihood maximisation. Section~\ref{sec:kriging}
derives simple and ordinary kriging as the GP posterior mean. Finally,
Section~\ref{sec:RKHS} establishes the RKHS connection, introduces the
reduced-rank approximation, and outlines kernel Kalman filtering.

% ============================================================
\section{Gaussian Processes}
\label{sec:GP_definition}
% ============================================================

\subsection{Definition and Basic Properties}

A stochastic process generalises the notion of a random variable to a
\emph{random function}: it is a family $(f(u))_{u \in \mathcal{U}}$ of random
variables indexed by a set $\mathcal{U}$, which may represent time
$\mathbb{R}_+$, physical space $\mathbb{R}^d$, or any other domain.

\begin{definition}[Gaussian process]
	For $d_{GP}\in\mathbb{N}^*$, a stochastic process $(f(u))_{u\in\mathcal{U}}$ on $\mathbb{R}^{d_{GP}}$ is a GP if every finite linear combination of its values follows a Gaussian distribution: for all $n\in\mathbb{N}^*$, all $u_1,\dots,u_n\in\mathcal{U}$, and all $a\in\mathbb{R}^n$, there exists $m\in\mathbb{R}^{d_{GP}}$ and $P\in\mathcal{M}_{d_{GP},d_{GP}}(\mathbb{R})$ with $P\succeq0$ such that:
	\begin{equation}
		\sum_{i=1}^na_if(u_i)\sim\mathcal{N}(m,P)
		\label{eq:GP_definition}
	\end{equation}
\end{definition}

A GP is associated with a \emph{mean function}:
\begin{equation}
	m:\left\{\begin{aligned}
	\mathcal{U}&\rightarrow\mathbb{R}^{d_{GP}}\\
	u&\mapsto\mathbb{E}[f(u)]
	\end{aligned}\right.
	\label{eq:GP_mean}
\end{equation}
and a \emph{covariance kernel}:
\begin{equation}
	k:\left\{\begin{aligned}
	\mathcal{U}\times\mathcal{U}&\rightarrow\mathcal{M}_{d_{GP},d_{GP}}(\mathbb{R})\\
	u,v&\mapsto\mathbb{E}[(f(u)-m(u))(f(v)-m(v))^\top]
	\end{aligned}\right.
	\label{eq:GP_kernel}
\end{equation}
We write $f\sim\mathcal{GP}(m,k)$ for a GP with mean function $m$ and covariance kernel $k$. As shown below, this notation is justified: $m$ and $k$ entirely characterise the law of $f$.
x
\begin{proposition}[Finite-dimensional marginals]
	\label{prop:GP_marginal}
	For any $n\in\mathbb{N}^*$ and any $u_1,\dots,u_n\in\mathcal{U}$, the joint
	distribution of $F:=(f(u_1),\dots,f(u_n))$ is the multivariate Gaussian:
	\begin{equation}
		\begin{pmatrix} f(u_1) \\ \vdots \\ f(u_n) \end{pmatrix}
		\sim \mathcal{N}\!\left(
			\begin{pmatrix} m(u_1) \\ \vdots \\ m(u_n) \end{pmatrix},
			\begin{bmatrix}
				k(u_1, u_1) & \cdots & k(u_1, u_n) \\
				\vdots       & \ddots & \vdots       \\
				k(u_n, u_1) & \cdots & k(u_n, u_n)
			\end{bmatrix}
		\right),
		\label{eq:GP_marginal}
	\end{equation}
	where each block $k(u_i,u_j)\in\mathbb{R}^{d_{GP}\times d_{GP}}$. In
	particular, the law of $f$ is entirely determined by $m$ and $k$.
\end{proposition}

\begin{proof}
	By \eqref{eq:GP_definition}, every linear combination $a^\top F$ (with
	$a\in\mathbb{R}^n$, applied component-wise to the $\mathbb{R}^{d_{GP}}$-valued
	blocks $f(u_i)$) is Gaussian, so $F$ is itself a Gaussian vector of
	$\mathbb{R}^{nd_{GP}}$. A Gaussian vector is entirely determined by its mean
	and covariance matrix, so it suffices to compute $\mathbb{E}(F)$ and
	$\mathrm{Cov}(F)$.

	On one hand, by linearity of the expectation and \eqref{eq:GP_mean}:
	\begin{equation}
		\mathbb{E}(F)=\begin{pmatrix}\mathbb{E}(f(u_1))\\\vdots\\\mathbb{E}(f(u_n))\end{pmatrix}=\begin{pmatrix}m(u_1)\\\vdots\\m(u_n)\end{pmatrix}.
	\end{equation}

	On the other hand, $\mathrm{Cov}(F)\in\mathcal{M}_{nd_{GP},nd_{GP}}(\mathbb{R})$
	is the block matrix whose $(i,j)$ block is, by \eqref{eq:GP_kernel}:
	\begin{equation}
		\mathrm{Cov}(F)_{i,j}=\mathbb{E}\left((f(u_i)-m(u_i))(f(u_j)-m(u_j))^\top\right)=k(u_i,u_j).
	\end{equation}

	Hence $F\sim\mathcal{N}\left((m(u_i))_{i=1}^n,\,(k(u_i,u_j))_{i,j=1}^n\right)$,
	which is \eqref{eq:GP_marginal}. Since this holds for every finite
	collection $u_1,\dots,u_n$, the whole family of finite-dimensional
	marginals of $f$ --- and hence the law of $f$ --- depends only on $m$ and
	$k$.
\end{proof}

A GP is thus a consistent infinite-dimensional generalisation of the
multivariate Gaussian distribution.

Both parameters admit a direct interpretation.
The mean function $m$ encodes prior knowledge about the average behaviour of
$f$, while the kernel $k$ encodes spatial and inter-output correlations:
the diagonal blocks $k(u,u) \in \mathbb{R}^{d_{GP} \times d_{GP}}$ measure the
covariance of each output with itself, while the off-diagonal blocks
$k(u,v)$ for $u \neq v$ capture how outputs at different locations co-vary.
In the context of this thesis, $m$ can be interpreted as a coarse model of
the observation function, and $k$ as a measure of the uncertainty or spatial
variability around this model.

\subsection{Covariance Kernels}

While the mean function $m$ typically reflects application-specific prior
knowledge about $f$ --- such as a known constant offset or a parametric
trend model --- it is the choice of covariance kernel $k$ that fundamentally
determines the regularity, smoothness, and correlation structure of the GP
sample paths. For $k$ to be a valid covariance kernel, it must be
\emph{symmetric positive semi-definite} (SPSD): $k(u,v) = k(v,u)^\top$ and, for any
$n \in \mathbb{N}^*$, $u_1, \ldots, u_n \in \mathcal{U}$, and
$a_1, \ldots, a_n \in \mathbb{R}^{d_{GP}}$,
$\sum_{i,j} a_i^\top k(u_i, u_j)\, a_j \geq 0$. We introduce two
classical families of kernels used throughout this thesis.

\paragraph{Squared exponential kernel.}
The squared exponential (SE) kernel, also called the radial basis function
(RBF) kernel, is defined by:
\begin{equation}
    k_{\text{SE}}(x, x') = \sigma_f^2\,\exp\!\left(
        -\frac{1}{2}(x'-x)^\top M (x'-x)
    \right), \quad M = \operatorname{diag}(l)^{-2},
    \label{eq:SE_kernel}
\end{equation}
where $\sigma_f^2 > 0$ is the signal variance and $l \in (\mathbb{R}_+^*)^d$
is the vector of length-scales. The parameter $\sigma_f^2$ controls the
amplitude of the process, while $l_i$ sets the characteristic distance along
dimension $i$ beyond which values of $f$ become approximately uncorrelated.
The SE kernel produces infinitely differentiable sample paths and is the most
widely used kernel in practice.

A sum of SE kernels with different length-scales can capture both local
fluctuations and global trends simultaneously:
\begin{equation}
    k(x, x') = \sum_{j=1}^J \sigma_{f_j}^2\,\exp\!\left(
        -\frac{1}{2}(x'-x)^\top M_j (x'-x)
    \right).
    \label{eq:SE_mixture}
\end{equation}

\paragraph{Mat\'{e}rn kernels.}
The Mat\'{e}rn family of kernels provides a flexible class of covariance
functions parametrised by a smoothness parameter $\nu > 0$:
\begin{equation}
    k_{\nu}(x, x') = \frac{2^{1-\nu}}{\Gamma(\nu)}
    \left(\frac{\sqrt{2\nu}\,\|x - x'\|}{l}\right)^\nu
    K_\nu\!\left(\frac{\sqrt{2\nu}\,\|x - x'\|}{l}\right),
    \label{eq:Matern_kernel}
\end{equation}
where $K_\nu$ is the modified Bessel function of the second kind and $l > 0$
is the length-scale. For $\nu = 1/2$, $3/2$, and $5/2$, closed-form
expressions exist; sample paths of $k_\nu$ are $\lceil\nu\rceil - 1$ times
differentiable. In particular, $\nu = 3/2$ and $\nu = 5/2$ are common choices
in practice as they strike a balance between the roughness of $\nu = 1/2$ and
the excessive smoothness of the SE kernel ($\nu \to \infty$).

\subsection{Hyperparameter Estimation}
\label{sec:hyperparameters}

The kernel $k$ depends on hyperparameters $\theta$ (e.g.\ $\sigma_f^2$ and $l$
for the SE kernel) that must be estimated from the training data. The standard
approach is type-II maximum likelihood (also called \emph{evidence maximisation}),
which maximises the marginal likelihood of the observations:
\begin{equation}
    p(\tilde{Z} \mid \tilde{X}, \theta)
    = \frac{\exp\!\left(-\frac{1}{2}
      (\tilde{Z} - m\,\mathbf{1})^\top K_\theta^{-1}
      (\tilde{Z} - m\,\mathbf{1})\right)}
      {\sqrt{(2\pi)^{\tilde{n}}\det(K_\theta)}},
    \label{eq:marginal_likelihood}
\end{equation}
where $K_\theta \triangleq k_\theta(\tilde{X}, \tilde{X}) + \tilde{R}\,
I_{\tilde{n}}$. Taking the logarithm and assuming $m = 0$, the
log-marginal-likelihood is:
\begin{equation}
    \ell(\theta) \triangleq \log p(\tilde{Z} \mid \tilde{X}, \theta)
    = -\frac{1}{2}\left(
      \tilde{Z}^\top K_\theta^{-1} \tilde{Z}
      + \log\det(K_\theta)
    \right) - \frac{\tilde{n}}{2}\log(2\pi).
    \label{eq:log_marginal_likelihood}
\end{equation}
The necessary conditions $\partial\ell/\partial\theta_i = 0$ yield, using the
differential of the matrix inverse and Jacobi's formula:
\begin{equation}
    \operatorname{tr}\!\left[\left(
    (K_\theta^{-1}\tilde{Z})(K_\theta^{-1}\tilde{Z})^\top - K_\theta^{-1}
    \right)\frac{\partial K_\theta}{\partial\theta_i}\right] = 0,
    \quad \forall\,i.
    \label{eq:MLE_hyperparameters}
\end{equation}
Since \eqref{eq:log_marginal_likelihood} is generally non-convex in $\theta$,
\eqref{eq:MLE_hyperparameters} is solved numerically by gradient-based
optimisation.

% ============================================================
\section{Gaussian Process Regression}
\label{sec:kriging}
% ============================================================

Kriging, also known as Gaussian process regression (GPR), is the problem of
estimating the value of $f$ at an unobserved location $x \in \mathbb{R}^d$
given noisy observations at a set of training points. We consider a GP
$f \sim \mathcal{GP}(m, k) : \mathbb{R}^d \to \mathbb{R}$ with known mean
function $m$. Without loss of generality, we work with the centred process
$f_c \triangleq f - m \sim \mathcal{GP}(0, k)$.

\paragraph{Setting.} Let $\tilde{X} = [\tilde{x}_1, \ldots,
\tilde{x}_{\tilde{n}}]^\top \in \mathbb{R}^{\tilde{n} \times d}$ be a set of
$\tilde{n}$ training inputs, and let:
\begin{equation}
    \tilde{Z} \triangleq
    \begin{bmatrix} f(\tilde{x}_1) + \tilde\varepsilon_1 \\ \vdots \\
    f(\tilde{x}_{\tilde{n}}) + \tilde\varepsilon_{\tilde{n}} \end{bmatrix}
    \in \mathbb{R}^{\tilde{n}}
    \label{eq:training_data}
\end{equation}
be the corresponding noisy observations, where
$\tilde\varepsilon_i \sim \mathcal{N}(0, \tilde{R}_i)$ are independent
observation noises. We seek to estimate $z \triangleq f(x) + \varepsilon$ at a
test point $x \in \mathbb{R}^d$, where $\varepsilon \sim \mathcal{N}(0, R)$.
We further assume uniform noise variance $\tilde{R}_i = \tilde{R}$ for all $i$,
so that $\tilde{Z} \mid f \sim \mathcal{N}(f(\tilde{X}), \tilde{R}\,
I_{\tilde{n}})$.

\paragraph{Best linear unbiased estimator.} Simple kriging seeks the best
linear unbiased estimator (BLUE) of $z$, i.e.\ the estimator of the form
$z^* = \Lambda^\top \tilde{Z}$ that is unbiased ($\mathbb{E}[z^* - z] = 0$)
and minimises the estimation variance $\operatorname{Cov}(z^* - z)$.

Unbiasedness is automatically satisfied since $\mathbb{E}[\tilde\varepsilon_i]
= \mathbb{E}[\varepsilon] = 0$. The estimation variance expands as:
\begin{equation}
    \operatorname{Cov}(z^* - z)
    = \Lambda^\top \underbrace{\left(k(\tilde{X}, \tilde{X})
      + \tilde{R}\,I_{\tilde{n}}\right)}_{G\,\triangleq}\Lambda
    - 2\Lambda^\top k(\tilde{X}, x) + k(x,x) + R,
    \label{eq:kriging_variance}
\end{equation}
where $G \in \mathbb{R}^{\tilde{n} \times \tilde{n}}$ is the \emph{Gram
matrix} of the kernel evaluated at the training points. Differentiating
\eqref{eq:kriging_variance} with respect to $\Lambda$ and setting to zero:
\begin{equation}
    \frac{d}{d\Lambda}\operatorname{Cov}(z^* - z) = 0
    \iff 2G\Lambda - 2k(\tilde{X}, x) = 0
    \iff \Lambda = G^{-1} k(\tilde{X}, x).
    \label{eq:kriging_weights}
\end{equation}

\subsection{Simple Kriging}
\label{sec:simple_kriging}

\begin{proposition}[Simple kriging estimator]
    \label{prop:simple_kriging}
    The simple kriging estimator and its posterior variance are:
    \begin{align}
        z^*_s(x, \tilde{X}, \tilde{Z})
        &= m + k(x, \tilde{X})\,G^{-1}\,(\tilde{Z} - m\,\mathbf{1}),
        \label{eq:simple_kriging_estimator} \\
        R^*_s(x, \tilde{X})
        &= k(x,x) + R - k(x, \tilde{X})\,G^{-1}\,k(\tilde{X}, x).
        \label{eq:simple_kriging_variance}
    \end{align}
\end{proposition}

\begin{proof}
    Substituting \eqref{eq:kriging_weights} into the expression for
    $z^* = m + \Lambda^\top(\tilde{Z} - m\,\mathbf{1})$ and into
    \eqref{eq:kriging_variance}:
    \begin{align}
        R^*_s &= \Lambda^\top G\Lambda - 2\Lambda^\top k(\tilde{X},x)
                 + k(x,x) + R \nonumber \\
               &= k(x,\tilde{X})\,G^{-1}\,G\,G^{-1}\,k(\tilde{X},x)
                  - 2k(x,\tilde{X})\,G^{-1}\,k(\tilde{X},x)
                  + k(x,x) + R \nonumber \\
               &= k(x,x) + R - k(x,\tilde{X})\,G^{-1}\,k(\tilde{X},x).
               \qedhere
    \end{align}
\end{proof}

\paragraph{Bayesian interpretation.} Since $z$ and $\tilde{Z}$ are jointly
Gaussian (both being noisy evaluations of the GP $f$), the simple kriging
estimator is also the posterior mean $\mathbb{E}[z \mid \tilde{Z}]$, and
$R^*_s$ is the posterior variance $\mathbb{V}[z \mid \tilde{Z}]$. To see
this, note that:
\begin{equation}
    \begin{pmatrix} z \\ \tilde{Z} \end{pmatrix}
    \sim \mathcal{N}\!\left(
        \begin{pmatrix} m \\ m\,\mathbf{1} \end{pmatrix},
        \begin{pmatrix}
            k(x,x) + R & k(x,\tilde{X}) \\
            k(\tilde{X},x) & G
        \end{pmatrix}
    \right),
    \label{eq:joint_GP}
\end{equation}
and applying the Gaussian conditioning formula \eqref{eq:Gaussian_conditioning}
directly yields \eqref{eq:simple_kriging_estimator}--\eqref{eq:simple_kriging_variance}.
Thus, kriging is the MMSE estimator of $z$ given $\tilde{Z}$ within the GP model.

\paragraph{Simultaneous estimation.} To estimate $Z = [f(x_1) + \varepsilon_1,
\ldots, f(x_n) + \varepsilon_n]^\top$ at $n$ test points simultaneously, the
scalar weight vector $\Lambda$ becomes a matrix $\Lambda \in \mathbb{R}^{\tilde{n}
\times n}$, and the estimator generalises to:
\begin{align}
    Z^*_s(X, \tilde{X}, \tilde{Z})
    &= m\,\mathbf{1} + k(X, \tilde{X})\,G^{-1}\,(\tilde{Z} - m\,\mathbf{1}),
    \label{eq:simple_kriging_batch} \\
    R^*_s(X, \tilde{X})
    &= k(X,X) + R\,I_n - k(X,\tilde{X})\,G^{-1}\,k(\tilde{X},X).
    \label{eq:simple_kriging_batch_var}
\end{align}
The diagonal entries of $R^*_s(X, \tilde{X})$ coincide with the individual
posterior variances at each test point.

\paragraph{Vector-valued outputs.} When $f : \mathbb{R}^d \to \mathbb{R}^{d_{GP}}$
with $d_{GP} > 1$, the kernel $k$ is matrix-valued and the Gram matrix becomes a
block matrix $G \in \mathbb{R}^{d_{GP}\tilde{n} \times d_{GP}\tilde{n}}$ with
$(i,j)$-block $k(\tilde{x}_i, \tilde{x}_j) \in \mathbb{R}^{d_{GP} \times d_{GP}}$.
The kriging estimator \eqref{eq:simple_kriging_estimator} and variance
\eqref{eq:simple_kriging_variance} remain formally identical, with scalar
products replaced by their matrix-valued counterparts. In this thesis, we
adopt the simpler assumption that the $d_{GP}$ output components are
\emph{a priori} independent, i.e.\ $k(u,v) = \mathrm{diag}(k_1(u,v), \ldots,
k_{d_{GP}}(u,v))$, which decouples the regression into $d_{GP}$ independent scalar
problems. This is a modelling choice, not a structural necessity of the GP
framework.

\subsection{Ordinary Kriging}
\label{sec:ordinary_kriging}

Simple kriging requires the mean function $m$ to be known. In many
applications, $m$ is constant but unknown. Ordinary kriging extends simple
kriging to this case by estimating $m$ from the data while maintaining the
BLUE property.

\paragraph{Setting.} The setup is identical to \cref{sec:simple_kriging},
except that $m \in \mathbb{R}$ is now an unknown constant. The unbiasedness
constraint $\mathbb{E}[z^* - z] = 0$ now becomes:
\begin{equation}
    \mathbb{E}[z^* - z] = 0
    \iff \left(\sum_{i=1}^{\tilde{n}} \lambda_i - 1\right) m = 0
    \iff \mathbf{1}^\top \Lambda = 1,
    \label{eq:OK_unbiasedness}
\end{equation}
since this must hold for any value of $m$.

\paragraph{Lagrangian formulation.} The constrained minimisation of
\eqref{eq:kriging_variance} subject to \eqref{eq:OK_unbiasedness} is solved
via Lagrange multipliers. Setting $\mathcal{L}(\Lambda, \mu) \triangleq
\operatorname{Cov}(z^* - z) + 2\mu(\mathbf{1}^\top\Lambda - 1)$ and
differentiating with respect to $\Lambda$:
\begin{equation}
    \frac{d\mathcal{L}}{d\Lambda} = 0
    \iff G\Lambda + \mu\,\mathbf{1} = k(\tilde{X}, x),
    \label{eq:OK_stationarity}
\end{equation}
which, combined with \eqref{eq:OK_unbiasedness}, gives the augmented linear
system:
\begin{equation}
    \begin{bmatrix} G & \mathbf{1} \\ \mathbf{1}^\top & 0 \end{bmatrix}
    \begin{bmatrix} \Lambda \\ \mu \end{bmatrix}
    =
    \begin{bmatrix} k(\tilde{X}, x) \\ 1 \end{bmatrix}.
    \label{eq:OK_system}
\end{equation}
Inverting \eqref{eq:OK_system} via the Schur complement of the upper-left
block $G$ (since the lower-right block $0$ is not invertible), with $S
\triangleq -\mathbf{1}^\top G^{-1} \mathbf{1} \in \mathbb{R}$:
\begin{equation}
    \Lambda = G^{-1}k(\tilde{X},x)
    + G^{-1}\mathbf{1}\;\underbrace{\frac{1 - \mathbf{1}^\top G^{-1}
    k(\tilde{X},x)}{\mathbf{1}^\top G^{-1}\mathbf{1}}}_{\in\,\mathbb{R}}.
    \label{eq:OK_weights}
\end{equation}

\begin{proposition}[Ordinary kriging estimator]
    \label{prop:ordinary_kriging}
    Defining the \emph{kriged mean} $m_o(\tilde{X}, \tilde{Z}) \triangleq
    \frac{\mathbf{1}^\top G^{-1} \tilde{Z}}{\mathbf{1}^\top G^{-1}
    \mathbf{1}}$, the ordinary kriging estimator and its posterior variance are:
    \begin{align}
        z^*_o(x, \tilde{X}, \tilde{Z})
        &= z^*_s(x, \tilde{X}, \tilde{Z},\, m_o(\tilde{X}, \tilde{Z})),
        \label{eq:OK_estimator} \\
        R^*_o(x, \tilde{X})
        &= R^*_s(x, \tilde{X}) + R^*_{m_o}(x, \tilde{X}),
        \label{eq:OK_variance}
    \end{align}
    where $R^*_{m_o}(x, \tilde{X}) \triangleq
    \dfrac{\left(1 - k(x,\tilde{X})\,G^{-1}\,\mathbf{1}\right)^2}
    {\mathbf{1}^\top G^{-1}\,\mathbf{1}} \geq 0$.
\end{proposition}

\begin{proof}
    Substituting \eqref{eq:OK_weights} into $z^* = \Lambda^\top \tilde{Z}$
    and rearranging yields \eqref{eq:OK_estimator}. For the variance, one
    computes $\Lambda^\top G \Lambda$ by substituting \eqref{eq:OK_weights},
    expanding the cross terms, and simplifying to obtain:
    \begin{equation}
        \Lambda^\top G \Lambda = k(x,\tilde{X})\,G^{-1}\,k(\tilde{X},x)
        + \frac{\left(1 - k(x,\tilde{X})\,G^{-1}\,\mathbf{1}\right)^2}
          {\mathbf{1}^\top G^{-1}\,\mathbf{1}},
    \end{equation}
    from which \eqref{eq:OK_variance} follows by substitution into
    \eqref{eq:kriging_variance}.
\end{proof}

Equation \eqref{eq:OK_estimator} shows that ordinary kriging is equivalent to
simple kriging with the unknown mean replaced by its data-driven estimate
$m_o$. The additional variance term $R^*_{m_o} \geq 0$ in \eqref{eq:OK_variance}
quantifies the extra uncertainty introduced by not knowing $m$. As the number
of training points grows, $m_o \to m$ and $R^*_{m_o} \to 0$, recovering simple
kriging asymptotically.

% ============================================================
\section{Kernel-Induced Functional Representations}
\label{sec:RKHS}
% ============================================================

Beyond regression, the covariance kernel of a GP induces further
functional representations exploited in this thesis: reduced-rank
approximations, which replace the full GP by a finite-dimensional projection
to reduce computational cost; the reproducing kernel Hilbert space, which
provides a deterministic function-analytic interpretation of GP regression;
and kernel Kalman filtering, which builds on this RKHS formulation to perform
recursive filtering directly in function space.

\subsection{Reduced-Rank Approximation}
\label{sec:reduced_rank}

The dominant computational cost of GP regression is the inversion of the Gram
matrix $G \in \mathbb{R}^{\tilde{n} \times \tilde{n}}$, which scales as
$\mathcal{O}(\tilde{n}^3)$. When $\tilde{n}$ is large, this is prohibitive.
Reduced-rank methods approximate the kernel by its truncated spectral
expansion.

Under mild conditions (Mercer's theorem), a symmetric positive semi-definite
kernel $k$ on a compact domain $\mathcal{X} \subset \mathbb{R}^d$ admits the
spectral decomposition:
\begin{equation}
    k(x, x') = \sum_{j=1}^\infty \mu_j\,\phi_j(x)\,\phi_j(x'),
    \label{eq:Mercer}
\end{equation}
where $(\mu_j, \phi_j)$ are the eigenvalue-eigenfunction pairs of the integral
operator $\mathcal{T}_k f \triangleq \int_{\mathcal{X}} k(\cdot, x) f(x)\,dx$,
ordered so that $\mu_1 \geq \mu_2 \geq \cdots \geq 0$. Truncating at rank $m
\ll \tilde{n}$ gives the approximation:
\begin{equation}
    k(x, x') \approx k_m(x, x') \triangleq \sum_{j=1}^m \mu_j\,\phi_j(x)\,\phi_j(x')
    = \Phi(x)^\top \operatorname{diag}(\mu_{1:m})\,\Phi(x'),
    \label{eq:reduced_rank_kernel}
\end{equation}
where $\Phi(x) \triangleq [\sqrt{\mu_1}\phi_1(x), \ldots, \sqrt{\mu_m}\phi_m(x)]^\top
\in \mathbb{R}^m$. Substituting \eqref{eq:reduced_rank_kernel} into the
kriging equations reduces the $\tilde{n} \times \tilde{n}$ inversion to an
$m \times m$ inversion via the Woodbury identity, bringing the cost down to
$\mathcal{O}(m^2\tilde{n})$.

\begin{remark}[Basis function interpretation]
    The reduced-rank GP is equivalent to a Bayesian linear model with feature
    map $\Phi$: the GP prior $f \sim \mathcal{GP}(0, k_m)$ corresponds to a
    weight-space prior $w \sim \mathcal{N}(0, I_m)$ with $f(x) = \Phi(x)^\top
    w$. This weight-space view is the starting point for random Fourier feature
    approximations and kernel recursive least-squares methods, outlined in
    Section~\ref{sec:kernel_kalman}.
\end{remark}

\subsection{Reproducing Kernel Hilbert Space Formulation}
\label{sec:RKHS_formulation}

\paragraph{Reproducing kernel Hilbert space.} Let $\mathcal{X} \subseteq
\mathbb{R}^d$ and let $k : \mathcal{X} \times \mathcal{X} \to \mathbb{R}$ be
a symmetric positive semi-definite kernel. The \emph{reproducing kernel Hilbert
space} $\mathcal{H}_k$ associated with $k$ is the completion of the linear
span $\{k(\cdot, x) : x \in \mathcal{X}\}$ under the inner product defined by:
\begin{equation}
    \langle k(\cdot, x),\, k(\cdot, x') \rangle_{\mathcal{H}_k}
    \triangleq k(x, x').
    \label{eq:RKHS_inner_product}
\end{equation}
The key property of $\mathcal{H}_k$ is the \emph{reproducing property}: for
any $f \in \mathcal{H}_k$ and $x \in \mathcal{X}$,
\begin{equation}
    f(x) = \langle f,\, k(\cdot, x) \rangle_{\mathcal{H}_k}.
    \label{eq:reproducing_property}
\end{equation}
This means that point evaluation is a bounded linear functional on
$\mathcal{H}_k$, which endows the space with strong regularity properties
determined by $k$.

\paragraph{Kriging as minimum-norm interpolation.} In the noise-free setting
($\tilde{R} = 0$), the simple kriging estimator \eqref{eq:simple_kriging_estimator}
coincides with the solution to the minimum-norm interpolation problem in
$\mathcal{H}_k$:
\begin{equation}
    z^*_s = \underset{f \in \mathcal{H}_k\,:\,f(\tilde{x}_i) = \tilde{z}_i}
    {\operatorname{argmin}}\;\|f\|^2_{\mathcal{H}_k}.
    \label{eq:RKHS_interpolation}
\end{equation}
In the noisy setting, replacing the interpolation constraint by a regularised
objective recovers \eqref{eq:simple_kriging_estimator} with $\tilde{R} > 0$
as the solution to:
\begin{equation}
    z^*_s = \underset{f \in \mathcal{H}_k}{\operatorname{argmin}}\;
    \left\{\|f\|^2_{\mathcal{H}_k}
    + \frac{1}{\tilde{R}}\sum_{i=1}^{\tilde{n}}
    (f(\tilde{x}_i) - \tilde{z}_i)^2\right\}.
    \label{eq:RKHS_regularisation}
\end{equation}

\paragraph{Correspondence between GP and RKHS.} The GP and RKHS perspectives
are complementary: the GP model provides a probabilistic interpretation of the
kriging estimator (posterior mean and variance under a Gaussian prior), while
the RKHS formulation provides a deterministic, functional-analytic
interpretation (minimum-norm interpolant in a function space). The two are
related through the Mercer representation \eqref{eq:Mercer}: the RKHS norm
penalises functions with large high-frequency components, with the decay rate
of $\mu_j$ controlling the smoothness enforced by $k$.

\begin{remark}[Implication for kernel choice]
    The equivalence \eqref{eq:RKHS_regularisation} shows that choosing a
    kernel amounts to choosing a notion of regularity for the unknown function
    $f$. A smooth kernel such as the SE kernel enforces infinitely
    differentiable interpolants, while a rougher Mat\'{e}rn kernel allows for
    functions with limited smoothness. In navigation applications, the choice
    of kernel should reflect prior knowledge about the physical field being
    modelled.
\end{remark}

\subsection{Kernel Kalman Filtering}
\label{sec:kernel_kalman}

\textbf{TODO.}
